\documentclass{pssbmac}

\usepackage[brazilian]{babel}

%\usepackage[latin1]{inputenc}
\usepackage[utf8]{inputenc} 

\usepackage[T1]{fontenc}
\usepackage{float}
\usepackage{graphics}
\usepackage{graphicx}
\usepackage{epsfig}
\usepackage{indentfirst}
\usepackage{amsmath, amsfonts, amssymb, amsthm}
\usepackage{url}
\usepackage{csquotes}

\newtheorem{theorem}{Theorem}[section]
\newtheorem{lemma}{Lemma}[section]
\newtheorem{proposition}{Proposition}[section]
\newtheorem{definition}{Definition}[section]
\newtheorem{remark}{Remark}[section]
\newtheorem{corollary}{Corollary}[section]
\newtheorem{teorema}{Teorema}[section]
\newtheorem{lema}{Lema}[section]
\newtheorem{prop}{Proposi\c{c}\~ao}[section]
\newtheorem{defi}{Defini\c{c}\~ao}[section]
\newtheorem{obs}{Observa\c{c}\~ao}[section]
\newtheorem{cor}{Corol\'ario}[section]

\usepackage[backend=biber, style=numeric-comp, maxnames=50]{biblatex}
\addbibresource{refs.bib}
\DeclareTextFontCommand{\emph}{\boldmath\bfseries}
\DefineBibliographyStrings{brazil}{phdthesis = {Tese de doutorado}}
\DefineBibliographyStrings{brazil}{mathesis = {Disserta\c{c}\~{a}o de mestrado}}
\DefineBibliographyStrings{english}{mathesis = {Master dissertation}}

\setcounter{topnumber}   {4}
\setcounter{bottomnumber}{4}
\setcounter{totalnumber}{10}
\DeclareMathOperator{\ulp}{ULP}

\renewcommand{\textfraction}     {0.15}
\renewcommand{\topfraction}      {0.85}
\renewcommand{\bottomfraction}   {0.70}
\renewcommand{\floatpagefraction}{0.66}

\setcounter   {dbltopnumber}           {2}
\renewcommand*{\dbltopfraction}        {1}
\renewcommand*{\dblfloatpagefraction}{0.9}

\usepackage{icomma}

\begin{document}

\title{Programação Funcional: do Cálculo Lambda à Implementação de Métodos Numéricos}

\author{
{\large Ribas, E. R. L. D.}\thanks{enzorochaeleitedinizribas@gmail.com}\, ,
{\large D'Afonseca, L. A.}\thanks{luis.dafonseca@gmail.com}\, ,
{\large Rocha, L. M.}\thanks{lucasmartinsrocha@gmail.com}  \\
{\small CEFET-MG, Belo Horizonte, MG} \\
}
\criartitulo
\vspace{-.5cm}
A programação funcional constitui um paradigma de programação fundamentado na utilização de funções como elementos centrais para a construção de programas. Diferentemente do paradigma imperativo, no qual algoritmos são frequentemente descritos por meio de alterações sucessivas no estado do programa, a abordagem funcional busca representar a computação por meio da composição e aplicação de funções, privilegiando conceitos como imutabilidade, funções de alta ordem, recursão e redução de efeitos colaterais. Uma das principais bases formais desse paradigma é o cálculo lambda, um sistema formal desenvolvido para representar funções e computações por meio de abstração, aplicação e substituição. Nesse contexto, expressões como $\lambda x.x$ permitem representar funções de maneira independente de uma linguagem de programação específica, estabelecendo uma relação entre a formulação matemática de uma função e sua representação computacional. A partir desses fundamentos, linguagens funcionais possibilitam expressar algoritmos de maneira próxima às suas formulações matemáticas, característica particularmente relevante no contexto de métodos numéricos. Neste trabalho, são introduzidos os principais conceitos associados ao paradigma funcional, com ênfase nos fundamentos do cálculo lambda e na relação entre funções matemáticas e construções computacionais. Como exemplo, considera-se o método de Newton-Raphson para determinação de raízes de funções, cuja formulação $x_{n+1}=x_n-\frac{f(x_n)}{f'(x_n)}$ pode ser interpretada sob uma perspectiva funcional, evidenciando a utilização de funções como abstrações fundamentais do algoritmo. A partir dessa abordagem, busca-se discutir as características da programação funcional e sua aplicabilidade à implementação de métodos numéricos, estabelecendo a fundamentação conceitual para a investigação comparativa entre diferentes paradigmas de programação.

% \nocite{*}
% \printbibliography

\end{document}

